% Source: https://artofproblemsolving.com/community/c3h1474577_quadratic_roots_from_vol_1
% Stable identity: aops-topic-1474577-q01
Let the equation be $p(t)=at^2+bt+c=0$ with rational $a\ne0$, $b$, $c$. For any rational $u$,
\[
p\!\left(x+u\sqrt z\right)=a\left(x^2+u^2z+2xu\sqrt z\right)+b\left(x+u\sqrt z\right)+c=P+Q u\sqrt z,
\]
where $P=a(x^2+u^2z)+bx+c$ and $Q=2ax+b$. Here $P$ depends on $u$ only through $u^2$, so with $u=y$ and $u=-y$:
\[
p(x+y\sqrt z)=P+Qy\sqrt z,\qquad p(x-y\sqrt z)=P-Qy\sqrt z ,
\]
with $P$ and $Qy$ rational.

We are told $P+Qy\sqrt z=0$. If $Qy\ne0$, then $\sqrt z=-\dfrac{P}{Qy}$ would be rational, a contradiction. So $Qy=0$, and then $P=0$ as well. Hence $p(x-y\sqrt z)=P-Qy\sqrt z=0$, so $x-y\sqrt z$ is a root. $\blacksquare$

(If $y=0$ the two numbers are equal and the claim is trivial; the argument covers this case too.)
